% Auto-generated by extract-marking.py -- do not edit by hand unless you know what to do.
% Works for BOTH \documentclass[...,type=solutions,...] and
% [...,type=marking,...] compiles, chosen at compile time via
% \OlympJuryPick / \OlympJuryBeginTable (see olympiad-marking.sty).
% \eqref{n} falls back to a plain '(n)' whenever the referenced
% equation isn't actually typeset in the current document.

\OlympJuryBeginTable{|p{0.88\linewidth}|>{\centering\arraybackslash}p{0.08\linewidth}|}{|p{0.66\linewidth}|>{\centering\arraybackslash}p{0.08\linewidth}|>{\centering\arraybackslash}p{0.06\linewidth}|>{\centering\arraybackslash}p{0.06\linewidth}|}
\hline
\OlympJuryPick{\multicolumn{1}{|c|}{\textbf{Содержание}} & \textbf{Баллы} \\ \hline}{\multicolumn{1}{|c|}{\textbf{Содержание}} & \textbf{Баллы} & \textbf{MP} & \textbf{A} \\ \hline}
\OlympJuryPick{Формула \eqref{1}: $\vec{v}_1(x,y,t) = \left(-\frac{v}{\sqrt{2}}, -\frac{v}{\sqrt{2}}\right) \cos\left(\omega t + \frac{\sqrt{2}\pi}{\lambda} x + \frac{\sqrt{2}\pi}{\lambda} y\right),$ & 0.4 \\ \hline}{Формула \eqref{1}: $\vec{v}_1(x,y,t) = \left(-\frac{v}{\sqrt{2}}, -\frac{v}{\sqrt{2}}\right) \cos\left(\omega t + \frac{\sqrt{2}\pi}{\lambda} x + \frac{\sqrt{2}\pi}{\lambda} y\right),$ & 0.4 & &  \\ \hline}
\OlympJuryPick{\quad\textit{\textbf{0.2 + 0.2 } за верную амплитуду вектора и зависимость фазы от координат} &  \\ \hline}{\quad\textit{\textbf{0.2 + 0.2 } за верную амплитуду вектора и зависимость фазы от координат} &  & &  \\ \hline}
\OlympJuryPick{Формула \eqref{2}: $\vec{v}_2(x,y,t) = \left(\frac{v}{\sqrt{2}}, -\frac{v}{\sqrt{2}}\right) \cos(\omega t - k_0 x + k_0 y),$ & 0.1 \\ \hline}{Формула \eqref{2}: $\vec{v}_2(x,y,t) = \left(\frac{v}{\sqrt{2}}, -\frac{v}{\sqrt{2}}\right) \cos(\omega t - k_0 x + k_0 y),$ & 0.1 & &  \\ \hline}
\OlympJuryPick{Формула \eqref{3}: $\vec{v}_3(x,y,t) = \left(-\frac{v}{\sqrt{2}}, \frac{v}{\sqrt{2}}\right) \cos(\omega t + k_0 x - k_0 y),$ & 0.1 \\ \hline}{Формула \eqref{3}: $\vec{v}_3(x,y,t) = \left(-\frac{v}{\sqrt{2}}, \frac{v}{\sqrt{2}}\right) \cos(\omega t + k_0 x - k_0 y),$ & 0.1 & &  \\ \hline}
\OlympJuryPick{Формула \eqref{4}: $\vec{v}_4(x,y,t) = \left(\frac{v}{\sqrt{2}}, \frac{v}{\sqrt{2}}\right) \cos(\omega t - k_0 x - k_0 y),$ & 0.1 \\ \hline}{Формула \eqref{4}: $\vec{v}_4(x,y,t) = \left(\frac{v}{\sqrt{2}}, \frac{v}{\sqrt{2}}\right) \cos(\omega t - k_0 x - k_0 y),$ & 0.1 & &  \\ \hline}
\OlympJuryPick{\quad\textit{Выше баллы \textbf{только} за верную амплитуду вектора \textbf{и} фазу} &  \\ \hline}{\quad\textit{Выше баллы \textbf{только} за верную амплитуду вектора \textbf{и} фазу} &  & &  \\ \hline}
\OlympJuryPick{Формула \eqref{5}: $v_y = 2\sqrt{2}v \cos(k_0 x) \sin(k_0 y) \sin(\omega t),$ , или эквивалентное уравнение двойной стоячей волны для хотя бы одной из компонент скорости & 0.8 \\ \hline}{Формула \eqref{5}: $v_y = 2\sqrt{2}v \cos(k_0 x) \sin(k_0 y) \sin(\omega t),$ , или эквивалентное уравнение двойной стоячей волны для хотя бы одной из компонент скорости & 0.8 & &  \\ \hline}
\OlympJuryPick{Формула \eqref{6}: $v_{\max} = \sqrt{8v^2} = 2\sqrt{2}v,$ & 0.5 \\ \hline}{Формула \eqref{6}: $v_{\max} = \sqrt{8v^2} = 2\sqrt{2}v,$ & 0.5 & &  \\ \hline}
\OlympJuryPick{Формула \eqref{7}: $\left( \frac{\lambda}{2 \sqrt{2}} \cdot n, \frac{\lambda}{2 \sqrt{2}} \cdot m \right) \ \ \text{где ровно  одно из \(n,m\) нечётное}.$ & 1.0 \\ \hline}{Формула \eqref{7}: $\left( \frac{\lambda}{2 \sqrt{2}} \cdot n, \frac{\lambda}{2 \sqrt{2}} \cdot m \right) \ \ \text{где ровно  одно из \(n,m\) нечётное}.$ & 1.0 & &  \\ \hline}
\OlympJuryPick{\quad\textit{\textbf{1.0} за полностью верный и обоснованный ответ для координат максимумов} &  \\ \hline}{\quad\textit{\textbf{1.0} за полностью верный и обоснованный ответ для координат максимумов} &  & &  \\ \hline}
\OlympJuryPick{\quad\textit{\textbf{0.4} за частично верный ответ, где видно что точки максимума будут на прямоугольной \textbf{бесконечной сетке}} &  \\ \hline}{\quad\textit{\textbf{0.4} за частично верный ответ, где видно что точки максимума будут на прямоугольной \textbf{бесконечной сетке}} &  & &  \\ \hline}
\OlympJuryPick{\textbf{Итого} & \textbf{3.0} \\ \hline}{\textbf{Итого} & \textbf{3.0} &  &  \\ \hline}
\end{tabularx}
